\section{Method}
\label{sec:method}

% AUTO-GENERATED by paper/regenerate.py -- do not hand-edit.
% id: fig:pipeline
% status: manual
% reason: manual block diagram -- author from method section. Source-of-truth: docs/CONCEPT.md.
\begin{figure}[t]
  \centering
  \fbox{\parbox{0.9\columnwidth}{\centering \textbf{fig:pipeline} -- placeholder.\\
  Status: manual.\\
  Gating step: ---}}
  \caption{DynamicMCPBench pipeline: server manifest $\to$ goal-gen $\to$ explorer $\to$ distiller $\to$ evaluator (Tier-1 / Tier-2) $\to$ report; the parallel `refresh` arrow for the living-bench loop.}
  \label{fig:pipeline}
\end{figure}


\subsection{Substrate and trace recording}
\label{sec:method:substrate}

We start from a \textbf{server manifest} (\texttt{dmcp/manifest.py})
declaring each MCP server's transport, \emph{dynamism class}
(\texttt{static} / \texttt{live\_read} / \texttt{stateful\_write}), and
a hard \texttt{sandbox=true} requirement for any
\texttt{stateful\_write} server. The recorder
(\texttt{dmcp/recorder.py}) is the only substrate the rest of the
pipeline talks to; its replay counterpart (\texttt{dmcp/replay.py}) is a
drop-in interchange that turns recorded traces into a deterministic
world for scoring (\S\ref{sec:method:refresh}).

\subsection{Forward generation: goal $\to$ exploration $\to$ distillation}
\label{sec:method:gen}

A \emph{goal generator} (\texttt{dmcp/goal\_gen.py}) drives an LLM to
propose realistic user goals against each server's tool surface
(per-server and cross-server pairs), with deterministic \emph{persona
seeding} (\texttt{dmcp/personas.py}). The \emph{explorer}
(\texttt{dmcp/explorer.py}) executes one goal as an agent loop with a
turn budget against the live substrate. On success, the recorded trace
is fed to the \emph{distiller} (\texttt{dmcp/distiller.py}), which
produces a \texttt{TaskSpec}:

\begin{itemize}
  \item a fuzzy natural-language \emph{prompt} (tool names stripped,
    concrete resources retained);
  \item one or more \emph{checkpoints}: \texttt{tool\_effect} (a tool
    from an equivalence set must have been called, optionally matching
    arg predicates) and \texttt{value\_produced} (a tool result must
    contain demanded substrings);
  \item optional \emph{minefields} (forbidden tools / args);
  \item an \emph{ordering} of checkpoints~--- partial; only where one
    effect genuinely depends on a prior one. Parallelizable effects are
    left unordered on purpose;
  \item a \texttt{ComplexityProfile} (trace depth, runtime branching,
    state coupling, cross-server) and dynamism class for stratification.
\end{itemize}

The distiller is explicit about its scope: it never invents checkpoints
the trace does not justify; it always echoes the LLM's ambiguity notes
into \texttt{TaskSpec.notes} so reviewers can see what the model was
unsure about.

% AUTO-GENERATED by paper/regenerate.py -- do not hand-edit.
% id: fig:trace_distill_example
% status: manual
% reason: manual worked example -- pick one v3 trace (preferably cross-server) and show its TaskSpec side-by-side.
\begin{figure}[t]
  \centering
  \fbox{\parbox{0.9\columnwidth}{\centering \textbf{fig:trace\_distill\_example} -- placeholder.\\
  Status: manual.\\
  Gating step: ---}}
  \caption{One worked example: a recorded trace (left) compiled into a `TaskSpec` (right) --- checkpoints, equivalence sets, minefields, partial order.}
  \label{fig:trace_distill_example}
\end{figure}


\subsection{Evaluation: Tier-1 deterministic, Tier-2 LLM judge}
\label{sec:method:eval}

Given a candidate trace and a \texttt{TaskSpec}, \textbf{Tier-1}
(\texttt{dmcp/evaluator.py}) is fully deterministic: for each
\texttt{tool\_effect}, verify some tool from the equivalence set was
called with arguments satisfying the predicate; for each
\texttt{value\_produced}, verify the substring; for each minefield,
verify it was not tripped; verify the partial order. \textbf{Tier-2}
(\texttt{dmcp/judge.py}) optionally upgrades failed
\texttt{tool\_effect} checks via an LLM \emph{effect-equivalence judge}
--- also constrained to a tool-call schema at temperature $0$. The fair
multi-agent path is \texttt{dmcp eval --replay}: candidates run against
a \texttt{TraceReplayRecorder} built from each spec's source trace, so
the world is identical across candidates and re-runs.

Critically, we \emph{never grade the final answer} for correctness. A
\texttt{value\_produced} checkpoint matches \emph{evidence the spec
demands}, never ``is the answer right''. Section~\ref{sec:rq1}~(RQ1)
shows what happens when one does.

\subsection{Living bench: refresh and decay}
\label{sec:method:refresh}

The refresh path (\texttt{dmcp/refresh.py}) re-runs each spec's
reference trace against the live substrate and classifies every
reference call as \emph{identical / drifted / broken / skipped}. A
\texttt{dmcp report} decay table turns this into the stratified drift
rate that signals when a spec should be re-distilled. Stateful-write
servers are \textbf{skipped by default} in refresh~--- re-running a
\texttt{git\_create\_branch} with the same name would just fail; the
explicit \texttt{--refresh-stateful} override is for sandboxes known
to support re-execution.

\subsection{Eval-side controls: tool-pool sampling, $P_{\text{alt}}$ degradation}
\label{sec:method:controls}

Beyond the headline scoring path, evaluation supports \textbf{distractor
pool modes} (\texttt{gold} / \texttt{target} / \texttt{full}) with a
sampler over six strategies (\texttt{random}, \texttt{hard\_neg},
\texttt{cross\_domain}, \texttt{same\_name}, \texttt{sibling},
\texttt{stratified}; \texttt{dmcp/sampling.py}) and a \emph{description
normalizer} (Level~A surface vs. Level~B semantic;
\texttt{dmcp/normalize.py}). The $P_{\text{alt}}$ driver
(\texttt{dmcp/curves.py}) sweeps the alternative-tool density and
produces degradation curves with Wilson confidence intervals by
complexity bin (\S\ref{sec:rq_palt}).
